% !TeX program = lualatex
% Compile twice: lualatex -interaction=nonstopmode -halt-on-error local_parity_decimation_note.tex
% Self-contained source. No external figures or bibliography files are required.
\documentclass[a4paper,11pt]{ltjsarticle}
\usepackage{luatexja-fontspec}
\setmainfont{Latin Modern Roman}
\setsansfont{Latin Modern Sans}
\setmonofont{Latin Modern Mono}[Scale=0.88]
\setmainjfont{HaranoAjiMincho-Regular.otf}[BoldFont=HaranoAjiMincho-Bold.otf]
\setsansjfont{HaranoAjiGothic-Medium.otf}[BoldFont=HaranoAjiGothic-Bold.otf]
\usepackage{amsmath,amssymb,amsthm,mathtools,bm,fix-cm}
\usepackage{geometry}
\geometry{top=23mm,bottom=24mm,left=24mm,right=24mm,headheight=15pt,headsep=7mm}
\usepackage{booktabs,array,longtable}
\usepackage{algorithm,algpseudocode,float}
\usepackage{xurl}
\usepackage[unicode,hidelinks]{hyperref}
\usepackage{fancyhdr}
\pagestyle{fancy}
\fancyhf{}
\fancyhead[L]{\small 局所パリティ周辺化による decimation 候補生成}
\fancyhead[R]{\small LPM-DP 1.0}
\fancyfoot[C]{\thepage}
\renewcommand{\headrulewidth}{0.3pt}
\setlength{\parskip}{0.3\baselineskip}
\setlength{\parindent}{1\zw}
\setlength{\emergencystretch}{2em}
\renewcommand{\arraystretch}{1.16}
\setcounter{tocdepth}{1}
\makeatletter
\renewcommand*\l@section[2]{\@dottedtocline{1}{0em}{4em}{#1}{#2}}
\makeatother
\numberwithin{equation}{section}
\theoremstyle{definition}
\newtheorem{definition}{定義}[section]
\newtheorem{proposition}[definition]{命題}
\newtheorem{remark}[definition]{注意}
\renewcommand{\proofname}{証明}
\floatname{algorithm}{アルゴリズム}
\newcommand{\Ftwo}{\mathbb F_2}
\newcommand{\N}{\mathcal N}
\newcommand{\V}{\mathcal V}
\newcommand{\A}{\mathcal A}
\newcommand{\I}{\mathbf 1}
\newcommand{\supp}{\operatorname{supp}}
\newcommand{\rank}{\operatorname{rank}}
\newcommand{\im}{\operatorname{Im}}
\newcommand{\clip}{\operatorname{clip}}
\newcommand{\TopK}{\operatorname{TopK}}
\newcommand{\softplus}{\operatorname{softplus}}
\newcommand{\lex}{\operatorname{lex}}
\newcommand{\code}[1]{\texttt{\detokenize{#1}}}
\newcommand{\ninf}{+\infty}
\DeclareMathOperator*{\argmin}{arg\,min}
\DeclareMathOperator*{\argmax}{arg\,max}
\hypersetup{pdftitle={局所パリティ周辺化による複数変数decimation候補の生成},pdfauthor={研究ノート},pdfsubject={LPM-DP 1.0: two-check marginal top-K candidate generator}}
\title{局所パリティ周辺化による\\複数変数 decimation 候補の生成\\[3mm]\large 2-check・最大4状態の動的計画法による実装仕様}
\author{研究ノート：LPM-DP 1.0}
\date{2026年9月23日}
\begin{document}
\maketitle
\thispagestyle{plain}
\begin{abstract}
BP が収束しなかったときに、次の BP に与える少数の hard-decimation パターンを生成する方法を定式化する。目的は完全な correction の探索ではなく、相互に関係する不確かな変数を、局所パリティ制約と整合する割当てで条件付けすることである。固定位置は LLR 履歴と check 間の共有変数から選び、固定値は最大2個の check を含む近似局所分布から選ぶ。固定しない近傍変数は最尤値に置換せず周辺化する。パリティを状態とする和積 DP と上位リスト DP を組み合わせ、固定変数数に対する全 $2^q$ 列挙を避ける。局所モデル内の top-$K$ の正確性、保持質量、計算量を示し、固定数の自動縮小、数値規約、warm start、および既存実装との接続を規定する。正確性は定義した局所モデルに対するものであり、全グラフの事後分布、BP の収束、論理誤り率に対する保証ではない。末尾では既存 SEARCH-BP-2.1 のスコアを検討する。
\end{abstract}
{\setlength{\parskip}{0pt}\tableofcontents}
\newpage

\section{何を改善し、何を保証する方法なのか}
\subsection{完全解の探索から条件付けの設計へ}
Beam Search Decoder は、不確かな1変数を $0/1$ に分岐し、その変数を message passing から除外して BP を継続する。各子の BP の反応を観測してから仮説を保持する点が、そのアルゴリズムの構造である\cite{beam}。本ノートは、この構造のうち候補生成だけを置き換える。固定後の BP と beam retention を不要にする方法ではない。

複数変数を一度に固定すると、複数の弱い入力や短い循環に同時に作用できる。しかし、独立に不確かな $q$ 変数に対して $K$ パターンだけを残すと、全割当てが等確率の場合の保持質量は $K/2^q$ である。固定数を増やすだけでは、正しい枝を捨てる危険が増える。本手法では、同じ check に結び付いた変数間のパリティ相関を使って値の組合せを評価し、少数候補で十分な局所質量を保持できないときは固定数を減らす。

Tesseract 型の $g+h$ は完全解の物理的コストの下界であり、有限反復 BP が成功する条件付けの良さを直接表す量ではない\cite{tesseract,dtd}。そこで本体には完全解探索、$f_{\mathrm{solve}}$、solve/guide の二系統 admission を含めない。以下では、変数の選択と値の選択を別の問題として定義する。

\subsection{既知の原理と本ノートの設計判断}
少数の parity-check 行を残し、その syndrome を trellis 状態として候補を生成する考え方は、古典符号の PC-GRAND に先行例がある\cite{pcgrand}。本ノートでは、チャネルに対する完全な誤り列ではなく、BP に渡す部分割当ての周辺確率を扱う。固定しない変数を和で消去するため、単純に joint configuration の上位候補を取り、その一部を抜き出す方法とも異なる。

check pair の選び方、固定位置の優先順位、保持質量による固定数制御、および既定パラメータは、本ノートで規定する設計判断である。これらを最適とする定理や、Beam Search Decoder に勝るという実測結果は与えない。一方、局所分布を定義した後の周辺化と top-$K$ 生成には、以下で示す正確性がある。この区別は、近似モデルの妥当性と実装の正しさを混同しないために必要である。

\section{入力、親 BP 状態、および出力の契約}
\subsection{復号問題と既存の固定}
$H\in\Ftwo^{m\times n}$ を復号用 parity-check 行列、$s\in\Ftwo^m$ を観測 syndrome とする。列は物理 qubit に限らず、DEM の fault variable でもよい。各列のチャネル確率 $p_j$ は $0<p_j\le1/2$ とし、物理的 LLR を
\begin{equation}
 w_j=\log\frac{1-p_j}{p_j}
     =\log1p(-p_j)-\log p_j
 \label{eq:channel}
\end{equation}
とする。対数は自然対数であり、正の LLR は $0$ を支持する。実際のノイズに相関があり独立 fault モデルを近似として用いている場合には、その近似も外側の復号問題の仮定である。

親 BP が既に固定している位置を $P$、値を $z_P$ とし、自由変数集合と残余 syndrome を
\begin{equation}
 V=\{0,\ldots,n-1\}\setminus P,\qquad
 \sigma=s\oplus H_Pz_P
 \label{eq:parent}
\end{equation}
と定める。既存の固定 $0$ も $1$ も $P$ に含める。現在の BP が解く問題は $H_Vx_V=\sigma$ である。$\N_V(a)=\{j\in V:H_{aj}=1\}$、$\N(j)=\{a:H_{aj}=1\}$ と書く。行・列の ID は原問題の0始まりの ID を最後まで使用する。固定パターンの辞書順は、昇順に並べた $(j,b_j)$ を平坦化した整数列の辞書順と定義する。

新しい候補は $D=(F,b_F)$、$\varnothing\ne F\subseteq V$ である。その候補に対する子問題は
\begin{equation}
 V_D=V\setminus F,\qquad
 \sigma_D=\sigma\oplus H_Fb_F,\qquad
 H_{V_D}x_{V_D}=\sigma_D.
 \label{eq:child}
\end{equation}
この式は正確な条件付けである。候補を採用すること自体は推測だが、採用後の行列と syndrome の縮約は近似ではない。

\subsection{必要な BP 情報}
入力は、親の固定マスク、最後まで完了した BP iteration の状態、直近 $W$ 回の posterior LLR 履歴である。履歴から使う量は
\begin{equation}
 \bar L_j=\frac1W\sum_{t=T-W+1}^{T}
           \clip(L_j^{(t)},-L_{\mathrm{hist}},L_{\mathrm{hist}}),
 \qquad j\in V,
 \label{eq:mean}
\end{equation}
である。clipping は平均の後ではなく各 iteration に先に施す。候補生成は少なくとも $W$ iteration を完了した非収束親に対してのみ呼び出す。短い履歴を暗黙に補完しない。

値の評価には、最後の完了 iteration で実際に使った check-to-variable message $\mu_{a\to j}^{(T)}$ を使用する。子の warm start には variable-to-check message $\nu_{j\to a}^{(T)}$ を使用する。この二つの方向を混同しない。後述する局所場を、履歴平均 $\bar L_j$ から check message を単純に引いて構成することも禁止する。時刻が異なるだけでなく、clipping や飽和後の減算は元の extrinsic sum を復元しないためである。

CSR と CSC の両方の隣接配列を持ち、各隣接 ID を昇順に保存する。重複辺は入力エラーとする。自由変数の LLR と message は有限でなければならず、NaN は直ちに拒否する。固定変数の $\pm\infty$ placeholder はすべてのスコア計算から除く。

\subsection{前処理と返却するもの}
入力時に、$\N_V(a)=\varnothing$ かつ $\sigma_a=1$ の行があれば親を矛盾として棄却する。$\sigma=0$ なら、全自由変数を $0$ とする correction が既に存在するので、外側の decoder が元の $H\hat x=s$ を確認して処理を終える。これは探索ではなく既存固定の検証である。自由変数が存在しても check に接続しない列は、位置選択の対象から外す。

正常終了時、生成器は同一の位置集合 $F$ を持つ高々 $K$ 個の異なる候補を、後述する局所コスト順で返す。各候補は親 ID、昇順の $(j,b_j)$ の組、局所負対数重み $c(D)$、正規化済み局所 log probability を持つ。候補群全体には選択 check の ID、実際の $q=|F|$、正規化定数、および切り捨て前の全局所分布を分母とする保持質量を付ける。異なる $F$ や異なる親の候補確率を足し合わせない。

\section{固定位置を選ぶ：不確かさと共有 check}
\subsection{不確かな変数の候補集合}
位置選択用の不確かさを
\begin{equation}
 u_j=1-\left|\tanh\frac{\bar L_j}{2}\right|
     =\frac{2}{1+\exp(|\bar L_j|)}
 \label{eq:uncertainty}
\end{equation}
と定める。大きな符号反転が履歴内で相殺される変数は $u_j$ が大きくなる。$u_j$ は確率の較正を主張しない順位付けの量である。単一時刻の $|L_j|$ より履歴を用いる動機は、Beam Search Decoder の summed-LLR 指標と共通する\cite{beam}。

check に接続する自由変数を $(-u_j,j)$ の辞書順に並べ、最初の $\min(M,|V_{\mathrm{active}}|)$ 個を $U$ とする。$M$ の既定値は32である。絶対閾値は設けず、全変数の $u_j$ が小さい場合にも相対的に不確かな変数を選ぶ。$U=\varnothing$ なら分岐せず終了する。

\subsection{曖昧な check の共有を明示的に計算する}
二つの異なる check $a<b$ に対して
\begin{equation}
 I_{ab}=\sum_{j\in U\cap\N_V(a)\cap\N_V(b)}u_j
 \label{eq:overlap}
\end{equation}
を定義する。上限 $r_{\max}=2$ のとき、$I_{ab}>0$ の pair があれば $(-I_{ab},a,b)$ の辞書順で最小の pair を選び、$\A=(a,b)$ とする。pair が存在しない、または $r_{\max}=1$ の場合には
\begin{equation}
 I_a=\sum_{j\in U\cap\N_V(a)}u_j
 \label{eq:single}
\end{equation}
が最大の check を選ぶ。同点は小さい行 ID を優先する。この場合 $\A=(a)$ である。正常な非空 $U$ からは必ず選択対象が得られる。

pair の生成は全 $m^2$ 組を走査しない。各 $j\in U$ について $\N(j)$ 内の異なる check pair に $u_j$ を加算する。変数を昇順に走査し、同じ pair の加算順を固定する。単一 check 用の $I_a$ も同じ走査で計算する。

この規則は $\sigma_a=1$ の行に限定しない。ここでの目的は syndrome の $1$ を消すことではなく、同じ不確かな変数によって結び付く推論問題を選ぶことである。共有変数が1個だけでも候補になり、そのとき4-cycle の存在までは意味しない。pair が2個以上の変数を共有すれば4-cycle が存在するが、本手法は cycle 数の最小化を解いているわけではない。

\subsection{局所近傍と固定優先順}
$\A$ の行数を $r\in\{1,2\}$ とし、局所近傍を
\begin{equation}
 V_\A=\bigcup_{a\in\A}\N_V(a)
 \label{eq:region}
\end{equation}
とする。位置候補は $U_\A=U\cap V_\A$ に制限する。各 $j\in U_\A$ の位置スコアを
\begin{equation}
 a_j=u_j\,|\N(j)\cap\A|
 \label{eq:position-score}
\end{equation}
とし、$(-a_j,j)$ の辞書順で並べる。その先頭
\begin{equation}
 q_0=\min(q_{\max},|U_\A|)
\end{equation}
個の順序列を $\pi=(j_1,\ldots,j_{q_0})$ とする。各 $1\le q\le q_0$ について
\begin{equation}
 F_q=\{j_1,\ldots,j_q\},\qquad B_q=V_\A\setminus F_q
 \label{eq:nested}
\end{equation}
を使う。$F_q$ は nested であり、固定数を減らすときに位置の再探索は行わない。

共有変数は同じ固定で二つの check に作用できるため係数2を与える。ただし、固定位置のスコアは値の正しさを判定しない。その判定に、次節の局所パリティ分布を使用する。

\section{固定値を評価する近似局所分布}
\subsection{選択した check の直接情報を除く}
選択した check を候補の評価に再度使うため、変数の事前重みには、その check 群からの直接 message を含めない。形式的な cavity field は
\begin{equation}
 \ell_j^{\mathrm{raw}}=w_j+
       \sum_{a\in\N(j)\setminus\A}\mu_{a\to j}^{(T)},
 \qquad j\in V_\A
 \label{eq:cavity-raw}
\end{equation}
である。これにより、同一の check 因子を unary なスコアと明示的な parity constraint の両方へ直接入れる重複を避ける。

実装上、上流 BP が非常に大きな有限値を使うことがあるため、本仕様では局所重みを作る段階だけで入力を有界化する。$\operatorname{cl}_L(x)=\min(L,\max(-L,x))$ と書き、規範的な局所場を
\begin{equation}
 \ell_j=\operatorname{cl}_{L_{\mathrm{prop}}}\!\left(
     \operatorname{cl}_{L_{\mathrm{prop}}}(w_j)+
     \sum_{a\in\N(j)\setminus\A}
           \operatorname{cl}_{L_{\mathrm{prop}}}(\mu_{a\to j}^{(T)})
                    \right)
 \label{eq:cavity}
\end{equation}
と定める。加算は check ID 順で行い、途中の累積和は clip せず最後に clip する。$L_{\mathrm{prop}}=30$ が既定値である。これは式\eqref{eq:cavity-raw}に加えた数値的な正則化であり、厳密な cavity posterior を計算したという主張ではない。BP 本体の message、$w_j$、および最終 correction の物理的コストは変更しない。

unary weight とその負対数を
\begin{align}
 \phi_j(b)&=\frac{1}{1+\exp((2b-1)\ell_j)},\qquad b\in\{0,1\},
 \label{eq:phi}\\
 e_j(b)&=-\log\phi_j(b)
       =\softplus((2b-1)\ell_j)
 \label{eq:unary-cost}
\end{align}
とする。すべての $\phi_j(b)$ は正であり、$\phi_j(0)+\phi_j(1)=1$ である。Min-sum BP の message から作る場合、この logistic 変換は局所モデルを定義する写像であって、較正された物理的確率への変換ではない。

\subsection{周辺化する対象を明確にする}
選択行の残余 syndrome を $\sigma_\A$ とし、$H_{\A,V_\A}$ を $H$ の対応する部分行列とする。局所分布は
\begin{equation}
 \widetilde P_\A(x_{V_\A})=
 \frac{1}{Z_\A}
 \left[\prod_{j\in V_\A}\phi_j(x_j)\right]
 \I[H_{\A,V_\A}x_{V_\A}=\sigma_\A]
 \label{eq:local-joint}
\end{equation}
である。選択外の check は hard constraint としては含めず、その情報の一部だけが局所場に入る。この分布は、元の $\Pr(x_V\mid H_Vx_V=\sigma,z_P)$ とは一般に異なる。

固定しない変数 $B_q$ の parity mass を
\begin{equation}
 Z_{B_q}(t)=
 \sum_{y\in\Ftwo^{|B_q|}:\ H_{\A,B_q}y=t}
       \prod_{j\in B_q}\phi_j(y_j)
 \label{eq:boundary-mass}
\end{equation}
とする。空集合では $Z_\varnothing(0)=1$、それ以外は0と定義する。すると固定パターンの周辺分布は
\begin{equation}
 \widetilde P_q(b)=\frac{1}{Z_\A}
       \left[\prod_{j\in F_q}\phi_j(b_j)\right]
       Z_{B_q}(\sigma_\A\oplus H_{\A,F_q}b).
 \label{eq:local-marginal}
\end{equation}
本手法の候補スコアは、この式の負対数である。

$B_q$ には、選択 check の全自由近傍から $F_q$ を除いた変数を入れる。$U$ の外側の変数を捨てたり、それらを $0$ に置換したりしない。また、$Z_{B_q}$ を各 check の別々の確率の積で置換しない。両 check に共有される $B_q$ の変数は、一つの変数として同時に二つのパリティへ寄与する。

\begin{proposition}[局所矛盾の意味]
$Z_\A=0$ であることと、$H_{\A,V_\A}x=\sigma_\A$ に解がないことは同値である。この場合、既存固定を維持した全体問題にも解がない。
\end{proposition}
\begin{proof}
すべての unary weight が正であるため、制約を満たす割当てが一つでもあれば正規化定数は正になる。逆に、全体解が存在すればその $V_\A$ への制限は選択行を満たす。
\end{proof}

\subsection{この局所モデルから期待する効果とその限界}
Min-sum の check 出力の絶対値は、宛先以外から受け取る入力の最小絶対値に制限される\cite{beam}。同じ check に弱い入力が複数ある場合、そのうち一つだけを固定しても別の弱い入力が残る。関連する複数変数を整合的に固定すると、残る自由変数への出力を制限する入力を同時に除ける場合がある。本手法は、この状況を少数の候補で扱うためのものである。

ただし、固定数を増やすと必ず message が正しい方向に強まるわけではない。外部 message にはループを通じた選択 check の情報が既に含まれ得る。近似局所モデルはその相関を再構成しない。固定後の BP も初期値に依存するため、同じ $D$ が同じ成功確率を持つという仮定は置かない。局所確率は候補生成の評価値であり、最後の成否は元の全グラフ上で判定する。

\section{最大4状態で周辺化と top-K を計算する}
\subsection{状態は変数の割当てではなくパリティである}
$\A=(a_0,\ldots,a_{r-1})$ は昇順とし、$r$ ビットのパリティを整数 $0,\ldots,S-1$、$S=2^r\le4$ で表す。変数ラベルと目標状態は
\begin{equation}
 h_j=\sum_{k=0}^{r-1}H_{a_kj}2^k,
 \qquad s_\A^{\mathrm{int}}=\sum_{k=0}^{r-1}\sigma_{a_k}2^k
 \label{eq:labels}
\end{equation}
である。以降の $t\oplus h_j$ は整数の bitwise XOR と解釈できる。行が線形従属でも最大4状態の配列をそのまま使い、到達不能な状態を無限コストとして扱う。行独立性を仮定しない。

計算を負対数領域で行うため、確率の加算に対応する演算を
\begin{equation}
 a\boxplus b=-\log(e^{-a}+e^{-b})
 =\min(a,b)-\log1p(e^{-|a-b|})
 \label{eq:logplus}
\end{equation}
と定義する。$a\boxplus\ninf=a$、$\ninf\boxplus\ninf=\ninf$ である。多項演算 $\mathop{\boxplus}_t$ は状態番号の昇順で左から畳み込む。確率の積は通常のコスト加算となる。$\boxplus$ と minimum を取り違えない。

一変数 $j$ を周辺化に追加する遷移は、任意の状態配列 $d$ に対して
\begin{equation}
 \operatorname{Step}_j(d)[t]
   =(d[t]+e_j(0))\boxplus(d[t\oplus h_j]+e_j(1))
 \label{eq:sum-transition}
\end{equation}
である。古い配列だけを読み、新しい配列へ書く二重バッファを使用する。

\subsection{固定しない変数の和積 DP}
まず $R=V_\A\setminus F_{q_0}$ を原変数 ID の昇順に処理する。初期配列は $d[0]=0$、$d[t]=\ninf$ $(t\ne0)$ とし、各 $j\in R$ に式\eqref{eq:sum-transition}を適用する。最終配列を $\mathsf D_{q_0}[t]$ とする。続いて $q=q_0-1,\ldots,0$ の順に
\begin{equation}
 \mathsf D_q=\operatorname{Step}_{j_{q+1}}(\mathsf D_{q+1})
 \label{eq:backward}
\end{equation}
を計算する。このとき
\begin{equation}
 \mathsf D_q[t]=-\log Z_{B_q}(t)
 \label{eq:backward-meaning}
\end{equation}
である。左辺の $\mathsf D_q[\cdot]$ は配列、右辺の添字 $B_q$ は式\eqref{eq:nested}の変数集合を意味する。実装では配列名を \code{suffix_cost[q][t]} として区別する。

固定数を減らすと、固定から外した変数は消えるのではなく $B_q$ 側へ戻る。式\eqref{eq:backward}は、この周辺化を毎回最初から計算せずに済ませるための再利用である。

\subsection{固定変数の全質量と上位リストを分けて保持する}
固定変数側にも全質量の配列 $A_q[t]$ を持つ。$A_0[0]=0$、$A_0[t]=\ninf$ $(t\ne0)$ から
\begin{equation}
 A_q=\operatorname{Step}_{j_q}(A_{q-1})
 \label{eq:forward}
\end{equation}
と進める。これは $F_q$ の全割当てを和で集約した値であって、上位 $K$ 候補だけの和ではない。

これとは別に、各 $q,t$ に上位リスト $\mathcal L_q(t)$ を保存する。リスト要素は $(c,\omega)$ であり、$c$ は $F_q$ の unary cost の和、$\omega$ は選択順 $\pi$ に沿った割当てのビット列である。初期状態は $\mathcal L_0(0)=\{(0,\varnothing)\}$、それ以外は空とする。

$\mathcal L_q(t)$ は次の二つの整列リストをマージし、先頭 $K$ 個を残すことで求める。
\begin{align}
 \mathcal E_0(t)&=
 \{(c+e_{j_q}(0),\omega0):(c,\omega)\in\mathcal L_{q-1}(t)\},
 \label{eq:extend-zero}\\
 \mathcal E_1(t)&=
 \{(c+e_{j_q}(1),\omega1):(c,\omega)\in\mathcal L_{q-1}(t\oplus h_{j_q})\},
 \label{eq:extend-one}\\
 \mathcal L_q(t)&=\TopK\bigl(\mathcal E_0(t)\cup\mathcal E_1(t)\bigr).
 \label{eq:list-dp}
\end{align}
順序は $c$ 昇順、厳密な同値なら選択順ビット列 $\omega$ の辞書順とする。$0<1$ である。同じ parity state に入った異なるビット列は別候補として保持する。state が同じだからという理由で一つにまとめてはいけない。

参照実装では $q_{\max}\le64$ とし、$\omega$ を unsigned 64-bit 整数として $\omega\leftarrow2\omega+b$ で更新する。同じ $q$ での整数順はビット列の辞書順と一致する。返却時には選択順と原 ID の対応を使ってビットを読み出し、$(j,b_j)$ を原 ID 順に並べ替える。$q=64$ では shift 幅64の式を使わず、最上位を含む各 bit を幅0--63の shift で取り出す。

\begin{proposition}[各状態で $K$ 個だけを保持してよい理由]
$\mathcal L_q(t)$ は、$F_q$ のうちパリティ $t$ を持つ全割当てから、unary cost とビット列の辞書順で上位 $K$ 個を正確に選んだリストである。
\end{proposition}
\begin{proof}
同一の親 state から同じ bit を付加する候補には、同じ cost が加算される。親で順位が $K$ より後ろの候補には、それより先に来る少なくとも $K$ 個の候補が存在し、それらも同じ遷移を利用できる。したがって、その候補は遷移先の上位 $K$ には入らない。初期状態からの帰納法で式\eqref{eq:list-dp}が成り立つ。cost が同じ場合にも、共通の末尾 bit の追加は辞書順を保存する。
\end{proof}

\subsection{周辺化済み境界重みを加えて最終候補にする}
固定数 $q$ ごとに、各 state のリストに境界 cost を加える。
\begin{equation}
 \mathcal C_q(t)=
 \left\{\left(c+\mathsf D_q[s_\A^{\mathrm{int}}\oplus t],\omega\right):
                    (c,\omega)\in\mathcal L_q(t)\right\}.
 \label{eq:terminal}
\end{equation}
境界 cost が $\ninf$ の state は丸ごと除く。高々4本の整列リストをマージし、全体の先頭 $K$ 個を $\mathcal C_q$ とする。この最終コストを $c_q(\omega)$ と書く。

\begin{proposition}[固定パターンの周辺 top-$K$]
$\mathcal C_q$ は式\eqref{eq:local-marginal}の確率が高い順の上位 $K$ パターンである。実数の厳密演算では、$q$ 変数の全 $2^q$ 列挙と同じ結果を得る。
\end{proposition}
\begin{proof}
式\eqref{eq:local-marginal}の非正規化確率の負対数は、固定変数の unary cost と、state だけに依存する境界 cost の和である。同じ state の全候補へ同じ定数を加えるため、その state 内で捨てた順位 $K$ より後ろの候補が全体の上位 $K$ に戻ることはない。したがって各 state の上位 $K$ を集めれば十分である。
\end{proof}

この手順では、$B_q$ の joint configuration の候補数を制限していない。$B_q$ は先に全て周辺化される。同じ $b_{F_q}$ に寄与する多数の境界割当ての質量を加算してから順位を付けることが必要である。joint configuration の top-$K$ を作った後で $F_q$ へ射影しても、一般には周辺 top-$K$ にはならない。

\subsection{正規化と保持質量}
正規化定数の負対数は、切り捨てを一切しない配列から
\begin{equation}
 \zeta_q=\mathop{\boxplus}_{t=0}^{S-1}
                  \bigl(A_q[t]+\mathsf D_q[s_\A^{\mathrm{int}}\oplus t]\bigr)
       =-\log Z_\A
 \label{eq:normalizer}
\end{equation}
と求める。実数演算では $q$ に依存しない。返却候補の局所 log probability は
\begin{equation}
 \log\widetilde P_q(\omega)=\zeta_q-c_q(\omega)
 \label{eq:normalized-cost}
\end{equation}
であり、その保持質量は
\begin{equation}
 C_K(q)=\sum_{\omega\in\mathcal C_q}\widetilde P_q(\omega),
 \qquad
 \log C_K(q)=\zeta_q-
                \mathop{\boxplus}_{\omega\in\mathcal C_q}c_q(\omega).
 \label{eq:coverage}
\end{equation}
候補リストだけを再正規化して合計1にしてしまうと、何を捨てたかが分からなくなる。固定数制御には必ず式\eqref{eq:normalizer}の全質量を分母として用いる。

\section{候補数を増やさず固定数を制御する}
\subsection{最大固定数から必要な分だけ縮小する}
$K\ge2$ を候補数の上限、$\tau\in(0,1]$ を局所保持質量の目標とする。採用固定数を
\begin{equation}
 q^*=\max\{q\in\{1,\ldots,q_0\}:C_K(q)\ge\tau\}
 \label{eq:q-selection}
\end{equation}
と定める。$q_0$ から1へ減らし、初めて条件を満たす $q$ の $\mathcal C_q$ を返す。$q=1$ の割当ては高々2個なので、局所モデルが可解なら $K\ge2$ により必ず全質量を保持できる。このため候補数を追加せずに終了する。

$K=2$、$q_{\max}=4$、$\tau=0.9$ を既定値とする。1親当たりの候補数は通常の binary branching と同じ2個のまま、局所相関または偏った局所重みがあるときにのみ複数固定を許す。これらの値は初期の参照仕様であり、性能が最適化された設定ではない。$q^*=1$ に戻っても、固定位置は共有 check に基づいて選んでいるため、元の Beam Search Decoder と同一の分岐になるとは限らない。

\begin{proposition}[保持質量の単調性]
$\A$、unary weight、および選択順 $\pi$ を固定したとき、実数演算で $C_K(q+1)\le C_K(q)$ である。
\end{proposition}
\begin{proof}
$F_{q+1}$ の任意の $K$ 個の割当てを $F_q$ へ射影すると、高々 $K$ 個の割当てになる。射影先の周辺確率の総和は元の総和以上である。その総和は、$F_q$ の top-$K$ の質量以下である。$F_{q+1}$ 側に top-$K$ を選べば主張を得る。
\end{proof}

固定数を変えるたびに局所場や選択 check まで更新すると、異なる局所モデルを比較することになり、この性質は失われる。本仕様では一回の候補生成中は親状態、$\A$、$\ell_j$、$\pi$ を変えない。

\subsection{2個の check で何がどこまで減るか}
局所重みが全て $\phi_j(0)=\phi_j(1)=1/2$ の場合を考える。$G_F=H_{\A,F_q}$、$G_B=H_{\A,B_q}$ とし、局所問題が可解であるとする。このとき固定値に残る独立な hard constraint の数は
\begin{equation}
 r_{\mathrm{eff}}=
 \rank[\,G_F\ G_B\,]-\rank(G_B)\le r\le2.
 \label{eq:effective-rank}
\end{equation}
非零の周辺確率を持つ $F_q$ の割当ては $2^{q-r_{\mathrm{eff}}}$ 個あり、すべて等確率なので
\begin{equation}
 C_K(q)=\min\left(1,\frac{K}{2^{q-r_{\mathrm{eff}}}}\right).
 \label{eq:uniform-coverage}
\end{equation}
これは候補欠落の制約を計算の工夫だけで取り除けないことを示す。

\begin{proof}
$G_Fb$ が $\sigma_\A+\im(G_B)$ に属するとき、かつそのときに限り境界変数で補完できる。商空間への写像 $b\mapsto G_Fb+\im(G_B)$ の rank は式\eqref{eq:effective-rank}であり、可解な fiber の大きさは $2^{q-r_{\mathrm{eff}}}$ である。各許容 $b$ には同じ個数の境界割当てがあるため周辺分布も一様になる。
\end{proof}

特に $G_B$ が選択行の全 rank を持ち、境界側が一様なら $r_{\mathrm{eff}}=0$ である。固定しない変数がどのパリティも吸収でき、選択 check は $F_q$ の割当て数を減らさない。default の $K=2$、$\tau=0.9$ ではこの状況で $q^*=1$ となる。一方、境界が空で独立2行が残る一様モデルなら、$q=3$ までを2候補で完全に保持できる。

非一様な境界重みでは、hard constraint の rank が減っても soft な相関が残り得る。その効果は $Z_{B_q}$ の state 間の違いとして正確に反映する。ただし、一般に2行だけで大きな曖昧領域を2候補へ圧縮できるわけではない。

\section{実装する処理を一意に定める疑似コード}
\subsection{局所領域の決定}
以下の疑似コードに現れる sort は、明記した全順序で実行する。式番号を参照する処理は、その式の入力と加算順をそのまま用いる。ハッシュ表の iteration 順序には依存させない。

\begin{algorithm}[H]
\caption{親状態から選択 check と固定優先順を作る}
\begin{algorithmic}[1]
\Require 親 $(H,s,P,z_P,\bar L)$、$M,q_{\max},r_{\max}$
\Ensure 昇順 check 列 $\A$ と優先順変数列 $\pi$
\State $V\gets\{0,\ldots,n-1\}\setminus P$、$\sigma\gets s\oplus H_Pz_P$
\State 前処理の矛盾・既存解を検査する（第2節）
\State 各 active な $j\in V$ の $u_j$ を式\eqref{eq:uncertainty}で計算する
\State $U\gets$ $(-u_j,j)$ 順の先頭 $M$ 個
\State 疎な表 $I_{ab}$ と $I_a$ を0で初期化する
\For{$j\in U$ を原 ID の昇順に走査}
  \For{$a\in\N(j)$}
    \State $I_a\gets I_a+u_j$
  \EndFor
  \For{$a<b$、$a,b\in\N(j)$}
    \State $I_{ab}\gets I_{ab}+u_j$
  \EndFor
\EndFor
\If{$r_{\max}=2$ かつ正の $I_{ab}$ が存在する}
  \State $\A\gets\argmin_{a<b}(-I_{ab},a,b)$
\Else
  \State $\A\gets(\argmin_a(-I_a,a))$
\EndIf
\State $V_\A\gets\bigcup_{a\in\A}\N_V(a)$、$U_\A\gets U\cap V_\A$
\State $U_\A$ を $(-u_j|\N(j)\cap\A|,j)$ 順に並べる
\State $q_0\gets\min(q_{\max},|U_\A|)$
\State \Return $\A$、この順序の先頭 $q_0$ 個からなる $\pi$
\end{algorithmic}
\end{algorithm}

\subsection{境界周辺化とリスト生成}
\begin{algorithm}[H]
\caption{全 $q$ に再利用できる state 配列と top-$K$ リストを作る}
\begin{algorithmic}[1]
\Require $\A,\pi,V_\A,K,L_{\mathrm{prop}}$、親の最後の check-to-variable message
\Ensure $\mathsf D_q,A_q,\mathcal L_q(t)$、$q=0,\ldots,q_0$
\State $S\gets2^{|\A|}$、$h_j,s_\A^{\mathrm{int}}$ を式\eqref{eq:labels}で計算する
\State 全 $j\in V_\A$ の $\ell_j,e_j(0),e_j(1)$ を式\eqref{eq:cavity}--\eqref{eq:unary-cost}で計算する
\State $d[0]\gets0$、$d[t]\gets\ninf$ $(t\ne0)$
\For{$j\in V_\A\setminus\{j_1,\ldots,j_{q_0}\}$ を原 ID の昇順に走査}
  \State $d\gets\operatorname{Step}_j(d)$
\EndFor
\State $\mathsf D_{q_0}\gets d$
\For{$q=q_0-1,\ldots,0$}
  \State $\mathsf D_q\gets\operatorname{Step}_{j_{q+1}}(\mathsf D_{q+1})$
\EndFor
\If{$\mathsf D_0[s_\A^{\mathrm{int}}]=\ninf$}
  \State \Return \code{LOCAL_INFEASIBLE}
\EndIf
\State $A_0[0]\gets0$、その他の $A_0[t]\gets\ninf$
\State $\mathcal L_0(0)\gets[(0,0)]$、その他のリストは空
\For{$q=1,\ldots,q_0$}
  \State $A_q\gets\operatorname{Step}_{j_q}(A_{q-1})$
  \For{$t=0,\ldots,S-1$}
    \State 式\eqref{eq:extend-zero},\eqref{eq:extend-one}の二つの整列リストを作る
    \State 和集合を $(c,\omega)$ 順で sort し、最初の $K$ 個を $\mathcal L_q(t)$ とする
  \EndFor
\EndFor
\State \Return 全配列とリスト
\end{algorithmic}
\end{algorithm}

\begin{algorithm}[H]
\caption{局所周辺 top-$K$ と保持質量から固定数を決める}
\begin{algorithmic}[1]
\Require アルゴリズム2の出力、$K\ge2$、$\tau\in(0,1]$
\Ensure 同一 $F_{q^*}$ を持つ候補群、$q^*$、保持質量
\For{$q=q_0,q_0-1,\ldots,1$}
  \State 各 $t$ のリストに $\mathsf D_q[s_\A^{\mathrm{int}}\oplus t]$ を加える
  \State 無限コストを除き、全 state の候補を sort した先頭 $K$ 個を $\mathcal C_q$ とする
  \State $\zeta_q\gets\mathop{\boxplus}_t(A_q[t]+\mathsf D_q[s_\A^{\mathrm{int}}\oplus t])$
  \State $\log C\gets\zeta_q-\mathop{\boxplus}_{(c,\omega)\in\mathcal C_q}c$
  \State 数値誤差の検査と $\log C$ の規約を第\ref{sec:numerics}節に従って適用する
  \If{$q=1$ または $\log C\ge\log\tau$}
    \State 各 $\omega$ を $\pi$ の先頭 $q$ 個への割当てへ変換する
    \State 各候補の $(j,b_j)$ を原 ID 順に並べる（候補順位は変更しない）
    \State 各候補に $\log\widetilde P=\zeta_q-c$ を付ける
    \State \Return $\mathcal C_q$、$q$、$C=\exp(\log C)$
  \EndIf
\EndFor
\end{algorithmic}
\end{algorithm}

候補生成中には BP を実行しないが、全 $q$ の値の組合せを列挙することもない。疑似コードの各 state リストは常に $K$ 以下である。$q$ を縮小するときは保存した prefix と suffix を組み合わせるだけなので、同じ大きな局所問題を何度も解き直さない。

\section{数値規約、例外処理、およびパラメータ}\label{sec:numerics}
\subsection{負対数領域で失ってはいけない情報}
参照実装の浮動小数点型は IEEE 754 binary64 とする。fast-math を無効にし、比較に epsilon を入れない。$\softplus(x)$ は
\begin{equation}
 \softplus(x)=\max(x,0)+\log1p(e^{-|x|})
 \label{eq:stable-softplus}
\end{equation}
で計算する。$\phi_j$ を先に丸めて $0$ や $1$ としてから対数を取らない。到達不能 state だけを $+\infty$ とし、非常に小さい確率を閾値で切らない。したがって、$Z_\A=0$ の判定は underflow ではなく到達可能性の判定となる。

式\eqref{eq:logplus}は両方が無限大の場合を分岐で処理する。無限大同士の減算、$0\cdot\infty$、NaN を計算に入れない。局所重みの入力 clipping は式\eqref{eq:cavity}だけで行い、DP cost に別の clipping を加えない。

数学上は整列リストのマージだけで足りるが、浮動小数点の定数加算は異なる cost を同値へ丸める場合がある。参照実装は各 state の高々 $2K$ 個の遷移候補を再度 $(c,\omega)$ で sort し、先頭 $K$ を取る。最終段階も高々 $SK$ 個を同じ順序で sort する。この規約は浮動小数点下の順位を一意にする。実数演算での top-$K$ 定理を、全割当ての浮動小数点列挙との bitwise 一致保証と解釈してはならない。

正規化の丸めにより $\log C_K(q)$ が微小に正となる場合を扱うため、$\epsilon_{\mathrm{mach}}=2^{-52}$ として
\begin{equation}
 \epsilon_{\mathrm{num}}=
 128\epsilon_{\mathrm{mach}}(1+|V_\A|+q_0)(1+L_{\mathrm{prop}})
 \label{eq:num-tolerance}
\end{equation}
を内部検査の許容量とする。$\log C>\epsilon_{\mathrm{num}}$ なら \code{NUMERICAL_ERROR} として処理を止め、それ以下の正値だけを0へ戻す。この閾値は誤差の厳密な上界ではなく、数値異常の検出規約である。候補同士の sort に使用しない。$q=1$ は全局所 support を保持するため、保持質量を解析的に1とする。それ以外での $\tau$ 判定は、規約適用後の $\log C\ge\log\tau$ とする。

\subsection{生成器のパラメータと影響}
生成器の自由度は表\ref{tab:params}の8個に限定する。その他の選択順、状態符号化、周辺化の範囲、および例外時の動作は固定仕様である。既定値には性能保証や測定による最適化の意味を持たせない。

\begin{table}[H]
\centering\small
\caption{候補生成器 LPM-DP 1.0 のパラメータ}
\label{tab:params}
\begin{tabular}{@{}p{34mm}p{16mm}p{27mm}p{70mm}@{}}
\toprule
名前 & 既定値 & 許容範囲 & 意味 \\
\midrule
\code{history_window} $W$ & 8 & 正整数 & posterior の符号振動を平均する窓幅。\\
\code{history_clip} $L_{\mathrm{hist}}$ & 25 & $0<L\le30$ & 履歴1点ごとの LLR 絶対値の上限。\\
\code{pool_size} $M$ & 32 & 正整数 & 位置選択に使う不確かな変数の上限。\\
\code{local_check_limit} $r_{\max}$ & 2 & $1$ または $2$ & 局所モデルに明示する check 数の上限。\\
\code{max_fixations} $q_{\max}$ & 4 & 整数 $1$--$64$ & 1回の候補で新たに固定する変数数の上限。\\
\code{candidates_per_parent} $K$ & 2 & 整数 $K\ge2$ & 同一親から返す局所周辺 top-$K$ の上限。\\
\code{retained_mass_target} $\tau$ & 0.9 & $0<\tau\le1$ & 固定数を許可する局所保持質量の目標。\\
\code{proposal_clip} $L_{\mathrm{prop}}$ & 30 & $0<L\le30$ & 局所場の入力と出力だけの clipping。\\
\bottomrule
\end{tabular}
\end{table}

$W$ を大きくすると長周期の符号反転を捉えやすくなる一方、最近の改善が過去の振動に埋もれやすい。$L_{\mathrm{hist}}$ を小さくすると極端な振幅の影響は減るが、信頼度の差も圧縮する。この二つは固定位置の選択に作用し、局所値の DP を直接変更しない。

$M$ は位置候補の探索範囲を広げる。大きくすると異なる check pair を選べるが、相対的に確かな変数も $U$ に入る。$r_{\max}=1$ では共有 check の推論を行わず2状態になり、$r_{\max}=2$ では共有変数の同時パリティを4状態で扱う。$r_{\max}>2$ は本仕様では拒否し、check 数に対する指数増加をパラメータ変更だけで導入しない。

$q_{\max}$ を大きくすると複数固定の機会は増えるが、実際の $q$ は保持質量で制限される。$K$ を大きくすると保持できる仮説数と外側の BP 回数が増える。$\tau$ を大きくするとより保守的な固定数になり、小さくすると局所質量の切り捨てを許容する。$K\ge2$ の制約は、$q=1$ への退避で両値を保持できるようにするためである。

$L_{\mathrm{prop}}$ は不確かな位置を選ぶ指標ではなく、局所分布の鋭さに関わる。小さすぎる値は境界の強い情報を弱める一方、大きい値は Min-sum の過信を局所モデルへ強く持ち込む可能性がある。局所保持質量が高くても、誤った境界情報を条件に高くなっている可能性は残る。

\subsection{例外と不正入力の処理}
行列形状、syndrome の二値性、固定マスクの値域、message 配列の長さ、および $W$ 点の履歴を検証する。$F_q$ と $P$ が重なること、同じ変数の重複指定、範囲外 ID、非有限パラメータは \code{INPUT_ERROR} とする。メモリ確保量と整数演算の overflow は呼出し前または確保前に検査し、黙って候補数を削らない。

選択局所問題が不整合なら \code{LOCAL_INFEASIBLE} を返し、その親を棄却する。別の check を選び直してこの矛盾を隠さない。返却候補が選択外の check と矛盾する可能性は別に残る。その判定は次節の hard-mask 適用時に行う。これは $Z_\A=0$ と異なり、局所モデルが見ていなかった制約との不整合である。

\section{hard masking、warm start、beam retention への接続}
\subsection{既存 snapshot に必要な追加情報}
参照した SEARCH-BP-2.1 の実装は、固定マスク、LLR 履歴、variable-to-check message を snapshot として保持し、子で free message を継承している\cite{repo}。本手法の式\eqref{eq:cavity}には最後の実際の check-to-variable message も必要である。既存 snapshot の説明にこの配列は含まれないため、候補生成用にその配列を公開・保存する拡張を明示的に行う。

最も単純な実装は $E=\operatorname{nnz}(H)$ 個の check-to-variable message を親 snapshot に追加する方法であり、binary64 なら親1個につき $8E$ byte 増える。選択領域を親状態の保存前に決めて必要な辺だけを保存する最適化も可能だが、最初の実装は全配列を保持する。$\nu^{(T)}$ から check 更新をやり直して作る message は時刻 $T+1$ の値であり、この仕様の $\mu^{(T)}$ と同一ではない。

\subsection{固定適用後の最初の BP 更新}
同じ親の候補は必ず同じ snapshot から独立に開始する。$D=(F,b_F)$ を適用するとき、式\eqref{eq:child}により全行の残余 syndrome を更新し、$F$ の値が $0$ でも $1$ でも全接続辺を BP の演算対象から除く。自由辺上の $\nu^{(T)}$ は親の値を保ち、チャネル prior は式\eqref{eq:channel}の元の $w_j$ のままとする。$\phi_j$、$\bar L_j$、局所 posterior を新しい prior にしない。

その後の最初の check 更新は、新しい residual と reduced neighborhood を使う。normalized min-sum を書けば
\begin{equation}
 \mu_{a\to j}^{\mathrm{new}}=
 \alpha(-1)^{(\sigma_D)_a}
 \prod_{k\in\N_{V_D}(a)\setminus\{j\}}
       \operatorname{sgn}(\nu_{k\to a}^{(T)})
 \min_{k\in\N_{V_D}(a)\setminus\{j\}}
       |\nu_{k\to a}^{(T)}|
 \label{eq:warm-update}
\end{equation}
である。その後で $w_j$ と新しい check message から posterior と次の variable-to-check message を計算する。親の check message は候補の評価には使ったが、子の check 更新を省略するためには使わない。

次数0で residual が1の check があれば、その候補は BP を回す前に棄却する。次数0で residual が0なら制約は自動的に満たされる。次数1の check は唯一の自由変数へその parity 値を要求する。参照接続では既存 kernel の次数1・飽和演算規約を維持し、候補生成器側で勝手に unit propagation を追加しない。符号と hard decision の零点規約も既存実装どおり、LLR が0なら bit 1 とする\cite{repo}。

子の履歴 ring は空にしてから再開する。親の過去の LLR を子の $W$ 点に混ぜない。少なくとも $W$ 回の更新後にも収束しなかった子だけを次の候補生成に渡す。元の $H\hat x=s$ を満たす子は直ちに有効 correction として扱い、正しい論理作用かどうかは実際の decoder 評価の対象とする。

\subsection{異なる親の局所スコアを直接比較しない}
$\widetilde P_q$ は親ごとに違う近似モデルに条件付けされている。さらに $q$ が違うと確率の対象事象も異なる。このため、全親の $-\log\widetilde P_q$ を一列に並べ、global top-$K$ を取る admission は採用しない。

参照接続は、最大 $B_{\mathrm{keep}}$ 個の各親から高々 $K$ 個を生成し、それら全てについて固定後の BP を評価する。$K=2$、$B_{\mathrm{keep}}=8$ なら1巡で最大16回である。これは既存の \code{k_run} を全親共通の上限として使う方式とは異なる。必要な処理数を途中で暗黙に8へ切り詰めない。

親の評価順は、前巡の BP retention 順とする。同じ親の候補は局所 cost と選択順ビット列の順で処理する。各子にこの順で一意の増加 ID を付け、最初の有効 correction で終了する。1巡で解が得られなければ、すべての非矛盾な子について BP 後のスコアを計算し、最大 $B_{\mathrm{keep}}$ 個を保持する。参照接続では既存の
\begin{equation}
 R(\text{child})=
 \frac1{|V_D|}\sum_{j\in V_D}
       \left|\tanh\frac{\bar L_j^{\mathrm{child}}}{2}\right|
 \label{eq:retention}
\end{equation}
を大きい順に使い、同点は全固定パターンの辞書順、さらに子 ID 順とする。異なる親から同じ固定パターンに至っても message が異なるので、固定パターンだけをキーに併合しない。

式\eqref{eq:retention}は BP 後に仮説を選別する既存の機構を維持するための参照規約であり、深さが異なる候補に対する最良のスコアと主張しない。複数固定で生じる比較上の偏りと代案は付録で扱う。

\begin{table}[H]
\centering\small
\caption{全体 decoder に接続する場合の参照設定。候補生成器自体の必須パラメータとは区別する。}
\begin{tabular}{@{}p{50mm}p{22mm}p{80mm}@{}}
\toprule
設定 & 参照値 & 作用 \\
\midrule
初期 BP iteration 上限 $T_0$ & 30 & root での推論時間。$T_0\ge W$ とする。\\
子 BP iteration 上限 $T_{\mathrm{child}}$ & 20 & 固定後の再推論時間。$T_{\mathrm{child}}\ge W$ とする。\\
保持親数 $B_{\mathrm{keep}}$ & 8 & 1巡後の BP 仮説数の上限。\\
最大巡数 $C_{\max}$ & 10 & 再帰的な候補生成の回数上限。\\
Min-sum scaling $\alpha$ & 1.0 & $0<\alpha\le1$。既存 kernel の設定をそのまま使う。\\
OSD fallback & false & 候補生成・BP のみの参照構成では無効。\\
\bottomrule
\end{tabular}
\end{table}

OSD を有効化する別構成では、既存実装の fallback を外側に残せる。ただし既存 OSD は固定値を必ず保つ制約付き solver ではなく、LLR の順位付けを使う別処理である\cite{repo}。その成否を本候補生成器の BP 収束と混同しない。最大巡数を使い切った場合、参照構成は宣言失敗を返す。新たな深さ探索、候補追加、または無制限の再始動を行わない。

\section{具体例：境界を消さずに候補を二つへ絞る}
局所 check を
\begin{equation}
 x_0\oplus x_1\oplus x_3=1,\qquad
 x_1\oplus x_2\oplus x_3=0
 \label{eq:example}
\end{equation}
とする。ここでは位置選択後の $F=\{0,1,2\}$ と $B=\{3\}$ を考える。固定側の unary weight は各 $1/2$、境界は $\phi_3(0)=0.9$、$\phi_3(1)=0.1$ とする。対応するラベルは
\begin{equation}
 h_0=1,\quad h_1=3,\quad h_2=2,\quad h_3=3,
 \qquad s_\A^{\mathrm{int}}=1
\end{equation}
である。境界の parity mass は $Z_B(0)=0.9$、$Z_B(3)=0.1$、その他は0となる。

$x_3=0$ で補完できる固定側の列は $011$ と $100$、$x_3=1$ で補完できる列は $001$ と $110$ である。正規化定数は $Z_\A=1/4$ なので、周辺確率は表\ref{tab:example}となる。

\begin{table}[H]
\centering
\caption{式\eqref{eq:example}の固定パターンの周辺分布。bit 列は $(x_0,x_1,x_2)$ 順。}
\label{tab:example}
\begin{tabular}{@{}cccc@{}}
\toprule
固定 bit 列 & 固定側 parity & 必要な境界 parity & $\widetilde P_F$ \\
\midrule
$011$ & $1$ & $0$ & $0.45$ \\
$100$ & $1$ & $0$ & $0.45$ \\
$001$ & $2$ & $3$ & $0.05$ \\
$110$ & $2$ & $3$ & $0.05$ \\
\bottomrule
\end{tabular}
\end{table}

$K=2$ なら最初の二つを返し、保持質量は $0.9$ である。したがって $\tau=0.9$ の実数仕様では3変数を二つの候補で固定できる。境界を最初から $x_3=0$ に固定する方法は、後半二つの質量を0にし、保持質量を誤って1と評価する。本手法は $x_3$ を実際には固定しないので、子 BP は $x_3$ と全ての領域外変数を引き続き推定する。

一方、$B$ が両パリティ方向の独立な一様変数を含む場合には $Z_B(t)=1/4$ が全 state で一定になり、固定側の全8列が等確率になる。このとき同じ3変数に対する二候補の質量は $1/4$ に過ぎず、固定数は縮小される。check の本数や共有変数数だけで圧縮効果を判断せず、境界周辺化後の質量を調べる必要がある。

\section{計算量、実装範囲、および保証の境界}
\subsection{候補生成は変数数に対して指数列挙しない}
最大変数次数を $d_v$、最大 check 次数を $d_c$、局所領域の全入射辺数を
\begin{equation}
 E_\A=\sum_{j\in V_\A}|\N(j)|
\end{equation}
とする。$U$ の選択はサイズ $M$ の heap を使って $O(n\log(M+1))$、共有 pair の生成は $O(Md_v^2)$ 個の加算である。決定的な連想配列を使う場合には pair 数に対する対数因子が付く。$|V_\A|\le r d_c$ であり、局所場の構成は $O(E_\A)$ である。

和積 DP と prefix/suffix の保存は $O(S|V_\A|)$ である。厳密演算の整列マージでのリスト DP と各 $q$ の候補選択は $O(SKq_0)$ である。binary64 参照実装の再 sort まで含めた安全な計算量表示は
\begin{equation}
 O\!\left(
 n\log(M+1)+Md_v^2\log(Md_v^2+2)
 +E_\A+S|V_\A|+SKq_0\log(SK+1)
 \right).
 \label{eq:complexity}
\end{equation}
ここで $S\le4$ である。境界変数数、固定数、符号長に対する $2^q$ や $2^{|B|}$ の項はない。

履歴平均が入力として既に存在する場合の式\eqref{eq:complexity}を示した。履歴から毎回計算するなら $O(nW)$ を加える。履歴 ring と running sum を親側に保持すれば更新は1 iteration 当たり $O(n)$ であるが、binary64 の加算順差が順位の境界に影響し得るため、実装の算術規約を version とともに固定する。

追加の局所作業メモリは、bit 列を64-bitで保持する参照仕様で
\begin{equation}
 O(Md_v^2+|V_\A|+SKq_0)
\end{equation}
である。これは親 snapshot や BP 自体の $O(E+nW)$ のメモリを含まない。既定値が固定された qLDPC family では候補生成は少なくとも全自由変数の評価走査を持ち、局所 DP の状態数だけが定数である。DP が4状態であることを decoder 全体が定数時間であるという意味に取らない。

\subsection{局所の正確さを全体の正確さへ拡張しない}
本手法が正確に扱うのは、式\eqref{eq:local-joint}で明示した最大2個の check と、式\eqref{eq:cavity}で定義した unary weight である。局所 top-$K$ はそのモデル内でのものであり、親固定が物理的に正しい確率、論理 sector の確率、BP が有限時間で収束する確率ではない。

選択外の check は同時に全て満たされる保証がない。そのため子の実行前には次数0の矛盾を調べ、BP 後には必ず元の $H\hat x=s$ を調べる。実行前検査で候補が減っても、別の候補を追加探索しない。生成時の保持質量はその削除前の局所質量であり、削除後や beam retention 後の全体被覆率ではない。異なる巡の保持質量を掛け合わせても、全体の正解保持確率にはならない。

複数変数の hard fixation は、その変数を通る cycle を実際に切断するが、選択規則が全グラフの cycle 数を最小にするわけではない。局所モデルの高い質量と loop の切断が、必ず正しい論理結果に繋がるわけでもない。改善の仮説は、個別 bit の判断では捉えにくい同一 check 近傍の相関を使い、保持する候補数を増やさずに関連する弱い入力をまとめて除去できる状況が存在する、という条件付きのものである。

\subsection{実装するモジュールの境界}
候補生成は \code{select_region}、\code{build_local_fields}、\code{build_dp_tables}、\code{choose_fixation_count} の4関数に分けられる。入力の親状態は全て immutable とする。境界 DP はメッセージを更新せず、子 BP は返却された位置と値だけを使って構造的に条件付けする。生成器内に OSD、Tesseract の frontier、全体 cost の下界、または BP callback を入れない。

本ノートの DP 式については、作成時に1--2行、1--8変数の600個の小規模入力を用い、可解な1,422個の固定集合について全列挙の周辺分布、top-$K$、正規化および保持質量の単調性と照合した。表\ref{tab:example}および境界が全 rank を吸収する一様例も照合した。これは有限例での演算検証であり、QEC の論理誤り率や復号時間を評価した数値実験ではない。

\appendix
\section{既存 SEARCH-BP-2.1 のスコアリングを見直す}
\subsection{参照した実装と付録の位置付け}
この付録は、\code{AKTKN/beamsearch_decimation} の branch
\code{hybrid-decoder-stage1}、commit
\path{9220771867197e797d4a9f9ccf4013cd13bcfe0f}
に含まれる \code{search_bp_scores.hpp}、\code{search_bp_decoder.hpp}、および実装説明を対象とする\cite{repo}。コード上の一般的性質を検討するものであり、別に保存された各実験の resolved parameter がこの commit の default と同じだったとは仮定しない。

付録の変更案は、本文の参照仕様へ暗黙に混ぜない。本文では候補生成を LPM-DP に置き換え、post-BP retention の機構を残した。そのうえで、残るスコアの解釈上の問題を整理する。

\subsection{下界の小ささと良い条件付けは異なる}
既存 \code{f_solve} は、固定値の物理 cost $g(D)$ と residual completion cost の fractional lower bound $h(D)$ の和である。これは
\begin{equation}
 g(D)+h(D)\le
 \min_{x:\ H_{V_D}x=\sigma_D}
      \left(g(D)+\sum_{j\in V_D}w_jx_j\right)
 \label{eq:old-bound}
\end{equation}
という方向の保証である。左辺が大きければ、その枝に安い解がないことを判断できる場合がある。しかし左辺が小さくても、実際に安い補完が存在するとは限らず、BP がそれを発見できるとも限らない。cost 下界を用いた探索の保証は、frontier を保った探索と正しい停止条件に関する保証である\cite{dtd}。

したがって、完全解探索を目的から外した本体では、この値を候補生成の primary score にしない。物理 cost は最終 correction を比較するために引き続き保持してよいが、局所 candidate posterior と加重平均して同一の意味を持つ量にしない。これが solve/guide の固定比率 admission を外す理由である。

\subsection{check の不確かさと不整合を区別する}
既存実装の check スコアは
\begin{equation}
 Q_a=\frac{1+v_a}{2},\qquad
 v_a=(-1)^{\sigma_a}
          \prod_{j\in\N_V(a)}\tanh(\bar L_j/2),
 \qquad \mathrm{ambiguity}_a=1-Q_a
 \label{eq:old-q}
\end{equation}
である\cite{repo}。$Q_a\simeq1/2$ は独立近似のパリティが不確かであることを表すが、$Q_a\simeq0$ はその近似が観測 parity と強く反対であることを表す。$1-Q_a$ はこの二つをまとめた不整合指標であり、確率的な uncertainty と同義ではない。

診断として両者を分けるなら、不確かさを $1-|v_a|$、不整合を $1-Q_a$ と別々に保存できる。hard decision の syndrome mismatch も別の量である。本体の位置選択は $Q_a$ を使わず、$u_j$ と実際の隣接関係により共有する不確かさを直接計算するので、この混同を避ける。

さらに、元の正しい条件付き分布 $\Pr(x\mid Hx=s)$ では全 check が確率1で満たされる。式\eqref{eq:old-q}は、その分布の真の check 成功確率ではない。既に相関している変数の周辺量を独立として掛け合わせた surrogate に過ぎない。本文の式\eqref{eq:local-joint}は、その代わりに選択 check の制約を局所モデル内で明示的に満たす。

\subsection{guidance gain の大きさがグラフサイズで消える}
既存実装は、$q=|F|$ に対して概略
\begin{align}
 J(D)&=\frac1q\sum_{j\in F}
       \softplus((2b_j-1)\bar L_j),\\
 G(D)&=\frac1q\left[
 \frac1{|V|}\sum_{j\in F}u_j+
 \frac\beta m\sum_{a\in\N(F)}(Q_a^D-Q_a)\right],\\
 f_{\mathrm{guide}}(D)&=J(D)-\lambda G(D)
 \label{eq:old-guide}
\end{align}
を用いている\cite{repo}。$\N(F)$ は固定に接する check の集合である。$0\le u_j\le1$、$|Q_a^D-Q_a|\le1$、$|\N(F)|\le qd_v$ から
\begin{equation}
 |G(D)|\le\frac1{|V|}+\frac{\beta d_v}{m}
 \label{eq:gain-bound}
\end{equation}
が従う。$\lambda,\beta,d_v$ が一定なら、グラフが大きくなるにつれ構造項は弱くなる。一方、$J(D)$ は全体変数数で小さくならない。したがって同じ \code{guidance_strength} が異なる問題サイズで同じ意味を持つとは限らない。

この surrogate を残す別案では、固定側は $q$、check 側は $|\N(F)|$ のような局所量で正規化すれば上のサイズ依存は除ける。ただし、既存 $G$ は固定後に BP を実際に更新した応答ではないため、正規化だけでは動的な収束を予測する問題は解消しない。本文では $J-\lambda G$ の微調整を行わず、局所周辺分布と保持質量へ置き換える。

$J$ を $q$ で割った値は joint assignment の負対数確率ではない点にも注意が必要である。本体の $c_q$ は割当ての unary cost を和として保持し、境界の周辺化項を加える。そのうえで、$q$ 間の優劣を raw cost では比較せず、同じモデルの保持質量で決める。

\subsection{post-BP retention が単なる削除を評価する偏り}
既存 retention の式\eqref{eq:retention}について、固定前の自由変数数を $n_V$、各変数の confidence を $c_j=1-u_j$、その平均を $R$ とする。$F$ を固定して削除しただけで、残りの posterior が全く変わらない場合でも
\begin{equation}
 R'-R=\frac{qR-\sum_{j\in F}c_j}{n_V-q}
 \label{eq:retention-bias}
\end{equation}
となる。平均より不確かな変数を消すと、BP が何も改善していなくても $R$ は上がる。1回に固定する数が一定の比較よりも、可変 $q$ の本手法ではこの効果を解釈する必要がある。

この偏りを除き、固定が残りの推論へ及ぼした効果だけを評価したい場合には、追加スカラー $d$ を親状態に保存する案がある。root で $d=0$ とし、新たに $F$ を固定した子に
\begin{equation}
 d_{\mathrm{child}}=d_{\mathrm{parent}}+
                    \sum_{j\in F}u_j^{\mathrm{parent}}
 \label{eq:debt}
\end{equation}
を渡す。retention の候補スコアを
\begin{equation}
 S_{\mathrm{prop}}(\mathrm{child})=
 \frac{d_{\mathrm{child}}+
         \sum_{j\in V_D}u_j^{\mathrm{child}}}{n}
 \label{eq:debiased-retention}
\end{equation}
とし、小さい順に保持する。分母 $n$ は decoder 全体で一定とする。同点は既存と同じ全固定パターン、子 ID の順を使う。この変更は1スカラーの継承と $O(n)$ の和で実装できる。

残りの $u_j$ が変わらなければ、削除した $u_j$ が $d$ に移るだけなので式\eqref{eq:debiased-retention}は変わらない。一方、固定が他の変数の不確かさを減らすとスコアは下がる。これは「消した変数の数」ではなく「固定後に残りへ伝わった情報」を測るための量である。

ただし、この代案は固定された変数の不確かさを履歴として残すため、多数の固定そのものが必要な正しい枝を不利にする可能性がある。真の事後確率や成功確率ではなく、単純削除の寄与を除いた動的な評価量である。したがって本文の参照 retention を無条件に置き換える仕様にはしない。有効 correction の発見は常にこのスコア比較より先に判定する。

\subsection{この付録から本体へ採用する変更}
本体で採用するのは、$f_{\mathrm{solve}}$ と $J-\lambda G$ による事前選別を止め、同一親内の局所周辺 top-$K$ を作ること、異なる親の局所 probability を直接競合させないこと、および固定後の BP による retention を残すことである。check uncertainty の分離と式\eqref{eq:debiased-retention}は独立の診断・変更案として扱う。複数の変更を同じスコアへ同時に足し込み、何を改善しようとしているのかを不明確にしない。

\begin{thebibliography}{9}
\small\setlength{\itemsep}{0.2em}\setlength{\parskip}{0pt}
\bibitem{beam}
M. Ye, D. Wecker, and N. Delfosse,
``Beam search decoder for quantum LDPC codes,''
arXiv:2512.07057v2 (2025).
特に Section II と Appendix B の masked BP、LLR 履歴、warm start を参照。
\url{https://arxiv.org/abs/2512.07057v2}.

\bibitem{pcgrand}
Y. Wang, J. Liang, and X. Ma,
``Partially Constrained GRAND of Linear Block Codes,''
arXiv:2308.14259v1 (2023).
少数の parity-check 行を trellis 制約に用いる先行例。本ノートの局所周辺 top-$K$ と保持質量制御そのものの出典ではない。
\url{https://arxiv.org/abs/2308.14259v1}.

\bibitem{tesseract}
L. Aghababaie Beni, O. Higgott, and N. Shutty,
``Tesseract: A Search-Based Decoder for Quantum Error Correction,''
arXiv:2503.10988v2 (2025).
\url{https://arxiv.org/abs/2503.10988v2}.

\bibitem{dtd}
K. R. Ott, B. Het\'enyi, and M. E. Beverland,
``Decision-tree decoders for general quantum LDPC codes,''
arXiv:2502.16408v1 (2025).
特に Sections 3.1, 5.1--5.2 の decimation と cost lower bound を参照。
\url{https://arxiv.org/abs/2502.16408v1}.

\bibitem{repo}
AKTKN, \code{beamsearch_decimation},
branch \code{hybrid-decoder-stage1}, commit
\path{9220771867197e797d4a9f9ccf4013cd13bcfe0f}.
2026年9月23日参照。対象ファイルは
\path{src/qec_bp_benchmark/native/search_bp_scores.hpp},
\path{src/qec_bp_benchmark/native/search_bp_decoder.hpp},
\path{docs/search_bp_implementation.md}。
\url{https://github.com/AKTKN/beamsearch_decimation/tree/9220771867197e797d4a9f9ccf4013cd13bcfe0f}.
\end{thebibliography}

\end{document}
